\begin{example}[Advanced 1: Height in a reducible affine scheme]\label{ex:vi31-adv-01}
For $A=k[x,y]/(xy)$, the prime $(x,y)$ has height $1$. The closed origin therefore has codimension $1$ even though two irreducible components pass through it.
\end{example}

\begin{example}[Advanced 2: Closed subset with components of different codimension]\label{ex:vi31-adv-02}
In $\mathbb A^3_k$, $W=V(x)\cup V(y,z)$ has components of codimensions $1$ and $2$. By convention $\operatorname{codim}(W,\mathbb A^3_k)=1$ and $\dim W=2$.
\end{example}

\begin{example}[Advanced 3: Principal ideal theorem]\label{ex:vi31-adv-03}
In a Noetherian domain, a minimal prime over a nonzero principal ideal has height exactly one. The nonzero condition excludes the generic component and the domain hypothesis makes height zero impossible.
\end{example}

\begin{example}[Advanced 4: Two equations need not give codimension two]\label{ex:vi31-adv-04}
$V(x,xy)\subset\mathbb A^2_k$ equals $V(x)$, so two displayed equations can define a codimension-one locus. Codimension counts prime chains, not equations.
\end{example}

\begin{example}[Advanced 5: Three equations with codimension two]\label{ex:vi31-adv-05}
$V(x,y,xy)\subset\mathbb A^3_k$ equals $V(x,y)$ and has codimension $2$, again showing redundancy among equations.
\end{example}

\begin{example}[Advanced 6: Product codimension]\label{ex:vi31-adv-06}
For coordinate subspaces $L_r\subset\mathbb A^m$ and $L_s\subset\mathbb A^n$ of codimensions $r,s$, their product has codimension $r+s$ in $\mathbb A^{m+n}$.
\end{example}

\begin{example}[Advanced 7: Graph of a morphism]\label{ex:vi31-adv-07}
For a morphism $f:X\to\mathbb A^m_k$ from a variety, its graph is isomorphic to $X$ inside $X\times\mathbb A^m_k$; when the graph is closed, its codimension is $m$.
\end{example}

\begin{example}[Advanced 8: Diagonal of a smooth-looking ambient product]\label{ex:vi31-adv-08}
For an $n$-dimensional separated variety $X$, the diagonal is closed in $X\times X$ and has codimension $n$ whenever the product has dimension $2n$.
\end{example}

\begin{example}[Advanced 9: Codimension and normalization]\label{ex:vi31-adv-09}
Let $A\subseteq B$ be an integral extension, as in normalization. Contraction of a strict chain of primes in $B$ remains strict by incomparability, while going-up lifts chains from $A$ once a starting prime upstairs is chosen. Thus integral extensions preserve Krull dimension globally, but equality $\operatorname{ht}(\mathfrak q)=\operatorname{ht}(\mathfrak q\cap A)$ for every specified prime $\mathfrak q$ is not a consequence of integrality alone; extra hypotheses such as going-down are needed for that stronger prime-by-prime statement.
\end{example}

\begin{example}[Advanced 10: Failure of global additivity from components]\label{ex:vi31-adv-10}
The scheme $\Spec k\sqcup\mathbb A^2_k$ has dimension $2$, but its point component has dimension and codimension both zero, so the additivity defect is $2$.
\end{example}

\begin{example}[Advanced 11: Nested additivity for varieties]\label{ex:vi31-adv-11}
If $Z\subset Y\subset X$ are subvarieties with dimensions $r<s<n$, then their codimensions are $s-r$, $n-s$, and $n-r$, so the relative codimensions add.
\end{example}

\begin{example}[Advanced 12: Codimension is not the number of generators]\label{ex:vi31-adv-12}
The ideal $(x^2,xy,xz)\subset k[x,y,z]$ has radical $(x)$, so its closed set has codimension $1$ although the ideal is displayed with three generators.
\end{example}

\begin{example}[Advanced 13: Generic-point minimum of local dimensions]\label{ex:vi31-adv-13}
For irreducible closed $Y$, the function $y\mapsto\dim\OO_{X,y}$ need not be constant, but its infimum on $Y$ is attained at the generic point and equals $\operatorname{codim}(Y,X)$.
\end{example}

\begin{example}[Advanced 14: Why catenary issues are deferred]\label{ex:vi31-adv-14}
Concatenating prime chains always gives a lower bound for total codimension, but equality for arbitrary Noetherian rings is not formal. Finite-type algebras over fields enjoy the dimension formulas needed in the variety case.
\end{example}